\documentclass[11pt]{article}

\usepackage[utf8]{inputenc}
\usepackage{geometry}
\usepackage{amsmath}
\usepackage{amsfonts}
\usepackage{amsthm}
\usepackage{amssymb}
\usepackage{hyperref}

\geometry{
	a4paper,
	margin=20mm
}

\usepackage{lastpage}
\usepackage{fancyhdr}
\pagestyle{fancy}

\usepackage[style=ieee, backend=biber, doi=false, isbn=false, eprint=false, url=false]{biblatex}
\addbibresource{bibliography.bib}
\appto\bibfont{\setlength{\emergencystretch}{.5em}}

\fancyhf{}

\allowdisplaybreaks

\title{
	Massively Parallel Domain Decomposition Method on a Modern GPU: Efficient Implementation and Fine-Tuning
}
% \author{Mieszko Grodzicki}
\date{}

\begin{document}

\maketitle

Exploiting GPU processing power for numerical computation requires regular data
structures and algorithms with a high degree of parallelism.
In this paper, we implement a massively parallel domain decomposition method \cite{Przykry}
and fine-tune the method's parameters for optimal efficiency on a~modern GPU.
Using the standard PyTorch framework with a small set of custom CUDA kernels
(such as a block sparse row matrix-vector product), we attain $70-90\%$ of the
measured peak memory bandwidth of an NVIDIA A100 GPU on individual solver
components. To further increase throughput, we leverage mixed-precision computation
in several parts of the method. In particular, we demonstrate that using half precision
in the subdomain solvers (for double-precision problems) yields substantial speedups
with negligible impact on convergence.

The resulting solver accepts the linear system explicitly as a sparse matrix, enabling
straightforward handling of variable or high-contrast coefficients and non-uniform meshes.
Despite the explicit input, the solver is memory-efficient: systems derived from
a~Discontinuous Galerkin discretization of elliptic problems with up to 100M degrees
of freedom (DoFs) in 2D and 50M DoFs in 3D can be solved on a single A100 with 80GB of memory.
The solver outperforms NVIDIA's AMG library, AmgX, in all tested scenarios,
with a maximum speedup of $4.92\times$ observed in 3D when high-order elements are used
for the discretization.

% TODO: format referencji taki jak chca w artykule?
\printbibliography

\end{document}
